\section{Two further projection-area identities}
\label{c:projection-appendix}

The barycentric argument in Section~\ref{c:rat-rational} gives the
reference area needed for the covering constructions.  We record two
independent geometric explanations of that area: one uses the second
moments of the face normals, and the other relates the projection of a
simplex to its longest parallel chord.

Use the edge-two tetrahedron
\[
 K_2=\{(x,y,ht):0\le t\le1,
                 \ |x|\le1-t,\ |y|\le t\},\qquad h=\sqrt2.
\]

\subsection{Second moments of the face normals}

\begin{lemma}\label{c:normal-frame}
For the regular tetrahedron $K_2$ just defined,
$A_{\min}(K_2)=\sqrt2$.
\end{lemma}
\begin{proof}
The facet-area form of the polytope projection formula
\cite[equation~(1), p.~83, and Section~3]{Martini1991} can be seen directly here.
The outward unit face normals are
\[
 n_1,n_2=(0,\pm\sqrt2,-1)/\sqrt3,
 \qquad n_3,n_4=(\pm\sqrt2,0,1)/\sqrt3.
\]
Each face has area $\sqrt3$, and for a unit vector $u$ the front faces
project onto the projection of $K_2$ with disjoint interiors, apart from
sets of area zero.  The back faces give the same area.  Consequently
\begin{equation}
 |\pi_{u^\perp}K_2|=\frac{\sqrt3}{2}\sum_{i=1}^4|u\cdot n_i|.
 \label{c:sing-normal-projection}
\end{equation}
Since $\sum_i n_i=0$ and $\sum_i n_i n_i^{\mathsf T}=\frac43 I_3$,
the numbers $a_i=u\cdot n_i$ sum to zero and have squared sum $4/3$.
If $S$ is the sum of their positive parts, then the negative parts have
absolute sum $S$, so $4/3=\sum_i a_i^2\le2S^2$.
Equation~\eqref{c:sing-normal-projection} is therefore at least
$\sqrt2$, with equality for $u=(1,0,0)$.  Thus $A_{\min}(K_2)=\sqrt2$.
\end{proof}

\subsection{Projection area from overlap volumes}

The same comparison area can be obtained from an identity for every
nondegenerate tetrahedron.  The simplex chord--projection identity
is attributed to Martini and Weissbach in
\cite[p.~8268]{Heinrich2014}.  The following proof differentiates the
volume of a translated intersection.

\begin{lemma}\label{c:nil-projection}
Let $T\subset\R^3$ be a nondegenerate tetrahedron, let $u$ be a
Euclidean unit vector, and let $L_u$ be the length of a longest chord of $T$
parallel to $u$.  Then
\begin{equation}\label{c:nil-chord-projection}
 |\pi_{u^\perp}T|=\frac{3\operatorname{vol}(T)}{L_u}.
\end{equation}
In particular,
$\min_{|u|=1}|\pi_{u^\perp}K_2|=h$.
\end{lemma}

\begin{proof}
Let $\lambda_1,\ldots,\lambda_4$ be the affine barycentric
coordinates of $T$, and put
$\Delta_i=\lambda_i(x+u)-\lambda_i(x)$.
These differences are independent of $x$ and sum to zero.  They do
not all vanish, since barycentric coordinates distinguish points.
Consequently
\[
 M=\sum_{i=1}^4\max(0,-\Delta_i)>0.
\]
For $e\ge0$, the conditions $x\in T$ and $x+eu\in T$ are
equivalent to
$\lambda_i(x)\ge e\max(0,-\Delta_i)$ for every $i$.
The lower bounds use total barycentric mass $eM$.  The intersection
$T\cap(T-eu)$ is therefore nonempty exactly when $e\le M^{-1}$;
for $e<M^{-1}$ it is a translate of a homothetic copy of $T$ of
ratio $1-eM$.  This proves $L_u=M^{-1}$ and
\begin{equation}\label{c:nil-overlap-volume}
 \operatorname{vol}(T\cap(T-eu))
 =(1-eM)^3\operatorname{vol}(T),\qquad 0\le e<M^{-1}.
\end{equation}

Let $\ell(q)$ be the length of the fiber of $T$ above
$q\in\pi_{u^\perp}T$.  Fubini's theorem gives
\[
 \operatorname{vol}(T\cap(T-eu))
 =\int_{\pi_{u^\perp}T}(\ell(q)-e)_+\,dq.
\]
Every point in the interior of the projection has a fiber of positive
length.  Indeed, homothetic copies of $T$ about an interior point,
with ratios increasing to one, lie in the interior of $T$, and their
projections exhaust the interior of its projection.  The boundary
of the projection polygon has area zero.  The difference quotient
of $(\ell-e)_+$ at $e=0$ lies between $-1$ and zero, and tends to
$-1$ wherever $\ell>0$.  Dominated convergence thus makes the right
derivative of the last integral equal to $-|\pi_{u^\perp}T|$.
Differentiating \eqref{c:nil-overlap-volume} proves
\eqref{c:nil-chord-projection}.

The horizontal sections of $K_2$ give
\[
 \operatorname{vol}(K_2)
 =h\int_0^1 4z(1-z)\,dz=\frac{2h}{3}.
\]
Its vertices $(\pm1,0,0)$ and $(0,\pm1,h)$ are pairwise at
distance two.  The distance between two convex combinations of the
vertices is at most the largest vertex distance, by the triangle
inequality.  Hence its diameter is two.  Every $L_u$ is at most
two, and an edge direction attains two.  The identity just proved
therefore gives the minimum projection area $h$.
\end{proof}
